\section{DOCSIS}
\label{sec:docsis}

DOCSIS (\textit{Data Over Cable Service Interface Specification}) es la familia de especificaciones de CableLabs para transmitir datos bidireccionales sobre redes de cable HFC \cite{docsis31, docsis40}, adoptada también por la UIT \cite{itu_j222}.

\subsection{Función dentro de la red de cable}
Una red de cable distribuye televisión en un solo sentido. DOCSIS le agrega un canal de datos bidireccional: reserva canales de RF del espectro descendente para los datos hacia el abonado y una banda baja para el sentido ascendente. Así, la misma red que lleva la TV ofrece acceso a Internet, telefonía IP y video sobre IP, sin un cableado nuevo.

\subsection{Equipos principales}
\begin{itemize}
    \item \underline{CMTS} (\textit{Cable Modem Termination System}): en la cabecera. Modula los canales descendentes, recibe los ascendentes y administra el acceso al medio: asigna a cada módem los intervalos en que puede transmitir.
    \item \underline{Red HFC:} fibra desde la cabecera hasta los nodos ópticos y, desde allí, coaxial con amplificadores bidireccionales, derivadores (\textit{taps}) y la acometida al abonado.
    \item \underline{Cable módem} (CM): en la casa del abonado. Sintoniza los canales DOCSIS, entrega la conexión por Ethernet o Wi-Fi y transmite en el sentido ascendente.
\end{itemize}

\begin{figure}[H]
    \centering
    \resizebox{\textwidth}{!}{%
    \begin{tikzpicture}[node distance=1.1cm, every node/.style={font=\small},
        caja/.style={draw, rounded corners, minimum height=1cm, minimum width=2.2cm, align=center, fill=blue!6}]
        \node[caja, fill=orange!12] (inet) {Internet /\\red IP};
        \node[caja, right=of inet] (cmts) {CMTS};
        \node[caja, below=0.6cm of cmts] (video) {Cabecera de\\video (QAM)};
        \node[caja, right=1.3cm of cmts] (comb) {Combinador\\RF};
        \node[caja, right=1.3cm of comb] (tx) {Transmisor y\\receptor óptico};
        \node[caja, right=1.6cm of tx] (nodo) {Nodo\\óptico};
        \node[caja, right=1.3cm of nodo] (amp) {Amplificador\\y \textit{taps}};
        \node[caja, above right=0.2cm and 1.0cm of amp, fill=green!8] (cm) {Cable\\módem};
        \node[caja, below right=0.2cm and 1.0cm of amp, fill=green!8] (stb) {Decodificador\\o TV};
        \draw[{Stealth}-{Stealth}] (inet) -- (cmts);
        \draw[{Stealth}-{Stealth}] (cmts) -- (comb);
        \draw[-{Stealth}] (video) -| (comb);
        \draw[{Stealth}-{Stealth}] (comb) -- (tx);
        \draw[{Stealth}-{Stealth}, very thick, red!70!black] (tx) -- node[above]{fibra} (nodo);
        \draw[{Stealth}-{Stealth}, very thick] (nodo) -- node[above]{coaxial} (amp);
        \draw[{Stealth}-{Stealth}] (amp.east) -- (cm.west);
        \draw[-{Stealth}] (amp.east) -- (stb.west);
        \node[draw, dashed, fit=(cmts)(video)(comb)(tx), inner sep=6pt, label=above:{Cabecera}] {};
    \end{tikzpicture}}
    \caption{Arquitectura de una red HFC con DOCSIS y televisión digital por cable.}
    \label{fig:docsis}
\end{figure}

\subsection{Sentidos de transmisión}
\begin{itemize}
    \item \underline{Descendente} (de la cabecera al abonado): el CMTS difunde a todos los módems del nodo. En DOCSIS 3.0 se usan canales QAM de $6 \, [MHz]$, de unos $38 \, [Mbps]$ cada uno, iguales a los de video digital. DOCSIS 3.1 introduce bloques OFDM de hasta $192 \, [MHz]$ con constelaciones de hasta 4096-QAM y codificación LDPC, y alcanza unos $10 \, [Gbps]$.
    \item \underline{Ascendente} (del abonado a la cabecera): el medio es compartido, así que los módems transmiten por turnos, en intervalos que asigna el CMTS (TDMA).
\end{itemize}

La \autoref{tab:docsis} resume la evolución de la familia.

\begin{table}[H]
    \centering
    \caption{Evolución de las versiones de DOCSIS (capacidades máximas aproximadas).}
    \label{tab:docsis}
    \begin{tabular}{c c c c p{5.4cm}}
        \toprule
        \textbf{Versión} & \textbf{Año} & \textbf{Descendente} & \textbf{Ascendente} & \textbf{Novedad principal} \\
        \midrule
        1.0 & 1997 & $40 \, [Mbps]$ & $10 \, [Mbps]$ & Primer estándar de módem por cable \\
        1.1 & 1999 & $40 \, [Mbps]$ & $10 \, [Mbps]$ & Calidad de servicio, telefonía IP \\
        2.0 & 2001 & $40 \, [Mbps]$ & $30 \, [Mbps]$ & Ascendente con A-TDMA y S-CDMA \\
        3.0 & 2006 & $1 \, [Gbps]$ & $200 \, [Mbps]$ & Agregación de canales, IPv6 \\
        3.1 & 2013 & $10 \, [Gbps]$ & $1$ a $2 \, [Gbps]$ & OFDM / OFDMA, LDPC, 4096-QAM \\
        4.0 & 2017 & $10 \, [Gbps]$ & $6 \, [Gbps]$ & Hasta $1{,}8 \, [GHz]$, \textit{Full Duplex} \\
        \bottomrule
    \end{tabular}
\end{table}

\subsection{DOCSIS, acceso a Internet y video sobre IP}
DOCSIS es, ante todo, la tecnología de acceso a Internet de las redes de cable: el CMTS funciona como un router de acceso y cada módem recibe una dirección IP. Como transporta IP, también sirve para distribuir video:
\begin{itemize}
    \item \underline{IPTV gestionada y OTT:} el decodificador recibe los programas como flujos IP por el módem. Es lo que hacen los decodificadores híbridos como el usado por Flow \cite{lightwave_flow}, que combinan la recepción QAM de los canales lineales con contenido IP a pedido.
    \item \underline{Migración a todo IP:} a medida que los operadores retiran los canales analógicos y QAM de video, ese espectro pasa a DOCSIS, y toda la televisión se entrega por IP.
\end{itemize}
